\chapter{Оптимизации физических операторов}

Предложенное решение по хранению JSON сохраняет преимущества
колоночной архитектуры только при условии, что физические операторы
исполнения способны эффективно использовать колоночное представление.
В ходе работы над реализацией обнаружился ряд узких мест, выявленных
профилированием на бенчмарке JSONBench: некоторые операторы,
унаследованные от документной подсистемы, сохраняли построчную
семантику, не используя в полной мере возможности SIMD и кэша~L1.
Настоящая глава подробно описывает сделанные оптимизации.

\section{Чтение только нужных колонок}

\subsection{Исходная проблема}

До оптимизации любой SELECT-запрос к колоночной таблице приводил
к чтению всех физических колонок таблицы, даже если в списке
SELECT и в WHERE упоминалось лишь несколько из них. Для таблицы с,
например, 50~колонками при запросе
\texttt{SELECT a, b FROM t WHERE c > 10;} с диска считывались все
50~колонок. Следствия:

\begin{enumerate}
    \item \textbf{Раздувание ввода-вывода.} Прочитанные, но
    неиспользуемые колонки впустую загружали страничный кэш,
    вытесняя полезные данные.
    \item \textbf{Лишняя аллокация и инициализация памяти.}
    Оператор \texttt{transfer\_scan} выделял буферы для всех
    колонок. При сканировании 10~млн строк из 50-колоночной таблицы
    это означало аллокацию сотен мегабайт никогда не используемой
    памяти.
\end{enumerate}

Для разреженных JSON-таблиц проблема обострялась многократно:
таблица с вычисляемой схемой может иметь десятки или сотни
физических колонок, и большинство запросов используют лишь одну-две
из них.

\subsection{Проход оптимизатора \texttt{column\_pruning}}

Реализовано новое правило оптимизации в
\texttt{components/planner/optimizer/rules/column\_pruning.cpp}
($\approx 350$~строк). Оно выполняется после этапа валидации
схемы (когда все логические ссылки на колонки уже разрешены в
физические индексы) и работает следующим образом:

\begin{enumerate}
    \item Обход логического плана в порядке BFS, поиск всех узлов
    \texttt{node\_aggregate\_t}. Каждый такой узел соответствует
    конечному вычислителю, обращающемуся к своей исходной таблице.
    \item Для каждого найденного \texttt{node\_aggregate\_t}~---
    сбор множества индексов колонок, фактически используемых
    запросом:
    \begin{itemize}
        \item Индексы, упомянутые в \texttt{node\_group\_t}
        (списке SELECT).
        \item Индексы, упомянутые в \texttt{node\_match\_t}
        (предложении WHERE).
        \item Для JOIN-запросов~--- колонки, участвующие в условии
        JOIN. Разделение по сторонам реализуется с учётом
        метаданных \texttt{key\_t::side()}.
    \end{itemize}
    \item Полученное множество индексов, отсортированное и без
    дубликатов, записывается в поле
    \texttt{node\_aggregate\_t::projected\_cols}.
    \item На этапе построения физического плана оператор
    \texttt{transfer\_scan} получает этот список через конструктор
    и использует его для создания <<разреженного>>
    \texttt{data\_chunk\_t}, в котором буферы выделяются только
    для запрошенных колонок.
\end{enumerate}

Совокупно это достигает двух целей одновременно: с диска читаются
только нужные колонки, и в памяти во время исполнения не
выделяются неиспользуемые буферы.

\subsection{Placeholder-векторы}

Введение placeholder-векторов потребовало аккуратной проверки всех
операторов на корректность обработки таких векторов. В
\texttt{components/vector/data\_chunk.cpp} добавлена вспомогательная
функция:

\begin{lstlisting}[language=C++,caption={Детекция placeholder-вектора},
                   label={lst:placeholder}]
static bool is_unprojected_placeholder(const vector_t& v) noexcept {
    return v.data() == nullptr && v.auxiliary() == nullptr;
}
\end{lstlisting}

Операции \texttt{data\_chunk\_t::copy},
\texttt{data\_chunk\_t::resize}, циклы итерации по колонкам во
всех операторах проверяют эту функцию и пропускают
placeholder-векторы. Индексы колонок остаются стабильными, но
физически никаких операций над непроекционными колонками не
производится.

\subsection{Измеренный эффект}

На бенчмарке JSONBench~Q1 (\texttt{SELECT COUNT(*) FROM bluesky}),
задействующем единственное поле, column~pruning сокращает время
запроса в 20--30~раз по сравнению с базовой версией. Эффект особенно
заметен при работе с полностью развёрнутой схемой (200+ физических
колонок). Для таблиц с вычисляемой схемой column~pruning критичен:
без него добавление новых JSON-путей делает все запросы всё
медленнее пропорционально росту числа колонок.

\section{Типизированные предикаты фильтрации}

\subsection{Проблема исходной реализации}

Оператор \texttt{operator\_match\_t}
(\texttt{components/physical\_plan/operators/operator\_match.cpp}) в
исходной версии выполнял фильтрацию построчно, с боксингом значений
в \texttt{logical\_value\_t}, виртуальными вызовами функций
сравнения и построчным копированием подошедших строк в выходной
чанк. На профилях JSONBench оператор \texttt{match} на запросах с
селективным WHERE занимал до 60\% времени выполнения.

\subsection{Скомпилированные типизированные предикаты}

Реализован механизм предварительной компиляции предиката в
типизированную функцию, выполняемую напрямую над буферами данных.
Логика работы~--- в
\texttt{components/physical\_plan/operators/predicates/predicate.cpp}:
при инициализации оператора \texttt{match} один раз на запрос
выражение предиката анализируется, определяются типы аргументов и
операций, создаётся функциональный объект, компилированный под
конкретные типы. Для предиката \texttt{col~INT > 10} создаётся
лямбда, работающая напрямую с \texttt{int64\_t*} буфером, без
боксинга и виртуальных вызовов. Этот предикат выполняется над
чанком пакетно: не <<строка-сравнение-строка>>, а
<<чанк-вычисление-bitmap>>.

\subsection{Пакетное копирование подошедших строк}

Вторая связанная оптимизация~--- переход от построчного копирования
к пакетному с использованием \texttt{indexing\_vector\_t}:

\begin{lstlisting}[language=C++,caption={Пакетное копирование через
                   indexing-вектор}, label={lst:batch-copy}]
std::pmr::vector<uint64_t> matched(resource);
matched.reserve(chunk.size());
for (size_t i = 0; i < chunk.size(); i++) {
    if (predicate->check(chunk, i)) {
        matched.push_back(i);
        if (!limit_.check(matched.size())) break;
    }
}
if (!matched.empty()) {
    vector::indexing_vector_t indexing(resource, matched.data());
    chunk.copy(out_chunk, indexing, matched.size(), 0);
}
\end{lstlisting}

Преимущества: один проход по каждой колонке вместо $O(n_\mathit{cols})$
проходов с боксингом; пакетный \texttt{vector\_ops::copy} с
indexing-вектором реализован как плотный цикл без виртуальных
вызовов; компилятор векторизует циклы пакетного копирования
автоматически для простых типов; ранний выход при LIMIT
применяется в чистом виде.

\subsection{Эффект}

Комбинация типизированных предикатов и пакетного копирования на
запросах JSONBench, содержащих селективный WHERE (например, Q3:
\texttt{\ldots WHERE commit.operation = 'create'}), ускоряет
фильтрацию в 5--8~раз. На запросах, где WHERE затрагивает
единственную числовую колонку и селективность предиката низкая,
выигрыш доходит до 15$\times$.

Столь значительный выигрыш объясняется тем, что исходная
реализация представляла собой почти наихудший сценарий для
современного процессора: плохая предсказуемость ветвлений (из-за
виртуальных вызовов), постоянные боксинг/анбоксинг на каждой
строке, отсутствие пакетной работы с кэш-линией.

\section{Ускорение GROUP BY}

\subsection{Профиль операторов агрегации}

Операторы агрегации в Otterbrix реализованы через пайплайн:
\texttt{operator\_group\_t} формирует хэш-таблицу, затем
\texttt{operator\_sort\_t} и \texttt{operator\_limit\_t}.
Профилирование показало, что \texttt{operator\_group\_t} при большом
числе уникальных ключей GROUP~BY становится узким местом.
Значительное время тратилось на вычисление хэша комбинированного
ключа, поиск с использованием универсальных compound-ключей, и
аккумуляцию через виртуальные вызовы агрегатных функций.

\subsection{Быстрый путь для простых случаев}

Значительная часть запросов JSONBench имеет форму
\texttt{SELECT f, COUNT(*) FROM t GROUP BY f ORDER BY COUNT(*) DESC LIMIT N;}
Для такого класса запросов введён специализированный быстрый путь:
вместо универсального compound-key-хэша используется хэш-функция,
оптимизированная под конкретный тип ключа; \texttt{COUNT(*)}
реализован без виртуальных вызовов; используется открытая адресация
хэш-таблицы, что даёт существенно лучшую локальность памяти по
сравнению с цепочечной адресацией. Совокупно быстрый путь ускоряет
запросы GROUP~BY~+~COUNT~ORDER~BY в 3--5$\times$ по сравнению с
универсальной реализацией.

\subsection{Исправление квадратичной сложности}

Отдельное улучшение связано с обнаружением квадратичной сложности в
одном из агрегатных операторов на запросах с большим числом
уникальных ключей. Вложенный цикл в реализации разрешения коллизий
имел сложность $O(n^2)$. Замена на линейное пробирование с
корректной обработкой deleted-tombstones довела сложность до
ожидаемой $O(n)$. На запросах JSONBench с 100~000 уникальных сессий
старый алгоритм завершался за десятки секунд, новый~--- за доли
секунды.

\section{Zone maps и обрезка row~groups}

\subsection{Принцип zone maps}

Zone~map~--- это мини-индекс, хранящийся вместе с каждым
row~group: для каждой колонки сохраняются агрегатные статистики
(min, max, count~null). При выполнении запроса с предикатом WHERE
оптимизатор имеет возможность пропускать целые row~groups, если их
zone~map заведомо исключает возможность совпадения. Для предиката
\texttt{WHERE commit.time\_us > 1700000000000} можно пропустить
все row~groups, в которых
\texttt{max(commit.time\_us) <= 1700000000000}.

\subsection{Обнаруженная ошибка и исправление}

В ходе работы обнаружилась ошибка в реализации zone~maps: при
работе с колонками с NULL-значениями zone~map иногда включал
некорректные min/max, вычисленные без учёта NULL-маски. В
результате row~groups, которые должны были пропускаться, не
пропускались, а иногда row~groups, которые должны были
обрабатываться, ошибочно пропускались (что приводило к неверным
результатам запроса). Исправление состояло в аккуратном учёте
NULL-маски при вычислении min/max: только ненулевые значения
участвуют в диапазоне.

\subsection{Взаимодействие с разреженными колонками}

Zone~maps применяются и к вспомогательным таблицам, хранящим
разреженные поля. Это даёт неочевидный, но важный эффект: предикат
по разреженной колонке может быть задействован для обрезки
row~groups соответствующей вспомогательной таблицы, а не всей
основной. Поскольку вспомогательные таблицы типично значительно
меньше основной, zone~maps по ним работают крайне эффективно.

\section{Оптимизация вставки}

Помимо оптимизаций чтения, существенные улучшения были сделаны на
стороне записи:

\subsection{Специализация \texttt{\_id} под int64}

Ранняя версия поддерживала для \texttt{\_id} несколько возможных
типов. Профилирование обнаружило, что обобщение приводило к боксингу
каждого идентификатора на входе и выходе MVCC-слоя. Специализация
\texttt{\_id} под единственный тип \texttt{int64\_t} привела к
двукратному ускорению пропускной способности \textsc{insert}.

\subsection{Параллельная запись в sparse-таблицы}

В первой версии реализации разреженного хранения каждая
вспомогательная таблица получала свой \textsc{insert} индивидуально.
На батче с 100~разреженными колонками это приводило к 100~отдельным
операциям подачи на шину актёров. Позднее операции были объединены
в параллельные батчи: все \textsc{insert} в разреженные таблицы
запускаются одной волной \texttt{co\_await} и ожидаются совокупным
\texttt{std::vector<future>}. Это уменьшило накладные расходы на
шине \texttt{actor-zeta} в 3--5$\times$.

\subsection{Пакетная вставка в jsonbench}

Для \texttt{jsonbench.cpp} внедрена пакетная вставка: вместо
\textsc{insert} по одной строке документы собираются в чанки
размером \texttt{INSERT\_BATCH~=~1000} и вставляются как один
\texttt{node\_insert\_t}.

\chapterconclusion

Описанные в настоящей главе оптимизации физических операторов, в
совокупности с архитектурой хранения из глав~3--4, обеспечивают
итоговую производительность системы, сопоставимую с промышленными
решениями на стандартных бенчмарках. Главные эффекты по категориям:
column~pruning сокращает ввод-вывод пропорционально избыточности
схемы; типизированные предикаты переводят оператор \texttt{match} в
пакетную парадигму; быстрые пути GROUP~BY и исправление
квадратичной сложности превращают агрегацию из узкого места в
дешёвую операцию; zone~maps с корректной обработкой NULL
обеспечивают грубую, но эффективную обрезку на уровне row~groups;
специализации и пакетизация \textsc{insert} обеспечивают высокую
пропускную способность записи.

Без этих оптимизаций теоретически правильная архитектура хранения
не дала бы практически конкурентоспособной производительности.
Именно их совокупность позволяет говорить о сопоставимых с
ClickHouse результатах, представленных в следующей главе.
